\chapter{Стационарность spin-branch determinant по масштабу}

Независимый аудит предложил последний маршрут без новой физики: использовать
scheme-independent разность periodic и anti-periodic fermion determinants и
проверить, фиксирует ли условие
\begin{equation}
 \frac{\partial}{\partial\chi}
 \Delta\Gamma_{\rm AP-P}(\chi)=0
\end{equation}
ненулевой scale. Одновременно аудит поставил вопрос о происхождении massive
vector. Оба пункта допускают точный ответ.

\section{Gauge mass не является внешним Stueckelberg input}

Finite block содержит charged coordinates
\begin{equation}
 x\mapsto e^{-i\theta}x,
 \qquad
 z\mapsto e^{+i\theta}z.
\end{equation}
В vacuum
\begin{equation}
 |x|=|z|=\frac{\chi}{\sqrt2}
\end{equation}
trace kinetic term даёт
\begin{equation}
 \Tr_{\rm orb}(D_\mu\Phi D^\mu\Phi)
 \supset4\chi^2A_\mu A^\mu.
\end{equation}
После canonical gauge normalization
\begin{equation}
 \frac{m_A^2}{\chi^2}=8g^2=3.
\end{equation}

Три scalar masses \(m_s^2=4\chi^2\) относятся к физическим координатам
quotient orbit space. Goldstone direction в этот список не входит: она
лежит вдоль gauge orbit и в \(R_\xi\)-комплексе имеет operator
\begin{equation}
 \Delta_0+\xi m_A^2,
\end{equation}
совпадающий с longitudinal/ghost block. Поэтому отсутствие massless scalar
среди трёх физических мод не означает отсутствия Higgs mechanism.

\begin{proposition}[Gauge-origin clarification]
Massive relative vector выводится из ковариантной kinetic metric charged
finite vacuum. Stueckelberg mass или уже проинтегрированный внешний сектор
для получения \(m_A^2=3\chi^2\) не требуются.
\end{proposition}

\section{Точная periodic/AP разность}

Положим
\begin{equation}
 x=\chi R_3,
 \qquad
 r=\frac{R_1}{R_3},
 \qquad
 a_k=k+\frac32,
 \qquad
 d_k=(k+1)(k+2).
\end{equation}
Для двух физических Dirac-пар
\begin{equation}
 \Delta\Gamma_{\rm AP-P}(x;r)
 =-2\sum_{k=0}^\infty d_k\,mathcal I(\rho_k),
 \qquad
 \rho_k=r\sqrt{a_k^2+x^2},
 \label{eq:spin-branch-mass-ratio}
\end{equation}
где
\begin{equation}
 \mathcal I(\rho)
 =\log\frac{\cosh(2\pi\rho)+1}{\cosh(2\pi\rho)-1}
 =2\log\coth(\pi\rho)>0.
\end{equation}
Все local heat-kernel coefficients двух branches совпадают. Поэтому в
\eqref{eq:spin-branch-mass-ratio} сокращаются \(\mu\), \(\lambda_0\),
\(\lambda_2\) и \(\lambda_4\).

\section{Аналитический derivative gate}

Поскольку
\begin{equation}
 \mathcal I'(\rho)
 =-\frac{4\pi}{\sinh(2\pi\rho)},
\end{equation}
получаем
\begin{equation}
 \frac{\partial\Delta\Gamma_{\rm AP-P}}{\partial x}
 =8\pi r x
 \sum_{k=0}^\infty
 \frac{d_k}
 {\sqrt{a_k^2+x^2}\,
  \sinh\left(2\pi r\sqrt{a_k^2+x^2}\right)}.
 \label{eq:spin-branch-positive-derivative}
\end{equation}
При \(r>0\) и \(x>0\) каждый член строго положителен. Следовательно,
\begin{equation}
 \frac{\partial\Delta\Gamma_{\rm AP-P}}{\partial x}>0
 \qquad(x>0).
\end{equation}
Единственная stationary point на замкнутой области \(x\ge0\) находится на
boundary \(x=0\). При \(x\to\infty\)
\begin{equation}
 \Delta\Gamma_{\rm AP-P}(x;r)\longrightarrow0^{-}.
\end{equation}

При \(r=1\) численный контроль даёт
\begin{center}
\begin{tabular}{ccc}
\toprule
\(x\) & \(\Delta\Gamma_{\rm AP-P}\) &
\(\partial_x\Delta\Gamma_{\rm AP-P}\) \\
\midrule
0 & \(-1.2984531\cdot10^{-3}\) & 0 \\
0.5 & \(-7.8082383\cdot10^{-4}\) & \(1.5474092\cdot10^{-3}\) \\
1 & \(-1.9482804\cdot10^{-4}\) & \(6.7652689\cdot10^{-4}\) \\
2 & \(-2.5002924\cdot10^{-6}\) & \(1.2468954\cdot10^{-5}\) \\
4 & \(-4.20\cdot10^{-11}\) & \(2.43\cdot10^{-10}\) \\
\bottomrule
\end{tabular}
\end{center}

\begin{theorem}[No nonzero scale from spin-branch ratio]
Periodic/AP determinant difference является конечным, scheme-independent и
нетривиально зависит от \(\chi R_3\), но не имеет stationary point при
\(\chi>0\). Поэтому она не фиксирует ни абсолютный scale, ни ненулевое
отношение \(\chi R_3\).
\end{theorem}

Тот же вывод не зависит от порядка вычитания: если определить обратную
разность \(\Gamma_{\rm P}-\Gamma_{\rm AP}\), derivative меняет знак, но
также нигде не обращается в нуль при \(x>0\).

\section{Итог внешнего аудита}

Независимый аудит верно подтвердил supertrace \(67\), determinant signs и
общий counterterm obstruction. Два уточнения имеют окончательный статус:
\begin{enumerate}
 \item gauge mass уже выведена из Higgsed relative-\(U(1)\) finite block;
 \item предложенный последний no-new-physics determinant route закрывается
 точной положительностью \eqref{eq:spin-branch-positive-derivative}.
\end{enumerate}

\begin{nogocriterion}[Закрытие relative-scale loophole]
Branch determinant считается scale selector только при наличии isolated
stationary point \(x_*>0\). Монотонная branch difference не может быть
использована для назначения \(\chi R\) по выбранному значению её magnitude.
\end{nogocriterion}

\begin{theorem}[Minimal fixed-\(K\) exhaustion]
После derivative gate минимальный fixed-\(K\) parent action исчерпан также
по последнему scheme-independent маршруту, не требующему новой структуры.
Переоткрытие требует либо новой chiral/topological measure, либо
variational carrier architecture, но не ещё одного determinant ratio внутри
текущей модели.
\end{theorem}

Машинный аудит сохранён в
\path{s2t_v4_spin_branch_mass_stationarity_gate_results.json}.